\subsection{What the class diagram currently decides}
\begin{proposition}[Proved: the class diagram with strict right inclusion]\label{D:known-diagram}
Among nonempty connected small categories,
\[
 \{\text{filtered}\}\subsetneq\{\text{exact}\}
 \subseteq\{\text{total}\}\subsetneq\{\text{quasi-filtered}\}.
\]
The ordinal-injection category witnesses the first strict inclusion.
The selected repaired AKR category of
Theorem~\ref{C:selected-nontotal} witnesses the last strict inclusion.
Whether the middle inclusion is an equality remains undecided here.
\end{proposition}
\begin{proof}
Filtered colimits of abelian groups are exact. The complex computing
$T_i(\Z[H])$ is the integral nerve-chain complex of $\El_AH$, so
exactness implies acyclicity of every component. We check the
homotopical step in supplied TOTASPH C.3--C.6. For a component
$E\subseteq\El_AH$ and $S:E\to\Set$, the functor
\[
 a\longmapsto\coprod_{(a,h)\in E}S(a,h)
\]
has category of elements isomorphic to $\El_E S$. Hence acyclicity
also holds for actions on $E$. For any cyclic subgroup
$C\leq\pi_1(BE,e)$, the action
$x\mapsto\Hom_{\Pi E}(e,x)/C$ gives the connected covering with
fundamental group $C$, as is checked by edge-path lifting. Its first
homology is $C$, forcing $C=0$. Thus $BE$ is simply connected,
and Hurewicz turns its acyclicity into weak contractibility. This
does not replace permutation modules by arbitrary diagrams.

For totality implying I, the checked TOTASPH C.11 argument first
gives common successors. If two principal cosieves were disjoint,
identify the values of the constant two-element action on their
union. Its element category is the double of $A$ identified on
those cosieves. Its reduced homology sequence has $H_1=\Z$,
since $BA$ is contractible and the two cosieves are nonempty and
connected. This contradicts totality. Applying common successors
to the inherited total category $p/A=\El_A A(p,-)$ completes
every span out of $p$ to a commuting square. The case $k=1$ of
Lemma~\ref{B:row-audit}, followed by common successors, gives
IV. Conditions I and IV give quasi-filteredness.

We checked the interpolation proof for $\Delta_{\mathrm{inj}}$
in supplied TOTASPH C.8--C.10. Its identity
$P_{t\delta^i}=y_iD_iP_t$ encodes gaps of injections as monomials.
The integral coordinates $x_i=z_i-z_{i+1}$ reduce extension of
compatible homogeneous restrictions to the alternating sum over
nonempty coordinate intersections. The prescriptions
$D_iC_{n,k}=C_{n-1,2k}\prod_{\ell\ne i}y_\ell^k$ agree on
intersections and have value one at the all-ones point. Thus they
give an augmentation-one combination equalizing the cofaces;
the simplicial identities then equalize all parallel injections.
Theorem 1 and Lemma 1 of \citet{IsbellMitchell1973}, or Theorem
3.1 and Proposition 3.3 of \citet{IsbellMitchell1976}, give
exactness. The two point inclusions $[0]\rightrightarrows[1]$
cannot be post-equalized by an injection. The 1973 Theorem 3
announces the example; the body of \citet{Isbell1974} was not
obtained and is not claimed as a proof checked in this run.

In repaired AKR, each parallel pair receives a quasi-filter and
each object pair a common successor, by the opened Constructions
B.9 and C.1 of \citet{AdamekKoubekReiterman1981}. These give I
and IV at the union: every finite diagram occurs at a finite stage.
Proposition~\ref{C:akr-all-rings} proves nonexactness for every
admissible repaired construction. Theorem~\ref{C:selected-nontotal}
gives a deterministic choice of sufficiently large targets for which
the binary category of elements has nonzero second homology. That
category is quasi-filtered and is not total, proving the final strict
inclusion without assuming the universal choice statement of C.1.
\end{proof}

\begin{obstruction}[Open in this incomplete record: Part D.1]\label{D:undecided-inclusions}
We have not decided whether exactness equals totality. The other
comparison is now decided: totality is strictly stronger than
quasi-filteredness, by Theorem~\ref{C:selected-nontotal}. The selected
AKR category is not total and therefore is not a separator for the
remaining middle inclusion.
\end{obstruction}

\subsection{The corrected manuscript statements}
\begin{proposition}[Refuted conjecture; known corrected implications]\label{D:manuscript}
The implication \emph{totalement $W_\infty$-asph\'erique
$\Longrightarrow$ filtrante} at II manuscript page 36ter is false.
For nonempty connected small $A$,
\[
 A\text{ filtered}\quad\Longleftrightarrow\quad
 A\text{ total and sifted}.
\]
For every fixed integer $n\geq1$,
\[
\begin{split}
 A\text{ filtered}\quad\Longleftrightarrow\quad&
 A\text{ total and }
 \hocolim_A K(M(a),n)\text{ is }n\text{-truncated}\\
 &\text{for every covariant }M:A\to\Ab.
\end{split}
\]
The implication from siftedness alone to totality in Proposition
2.19 at II manuscript page 36 is false.
\end{proposition}
\begin{proof}
Proposition~\ref{D:known-diagram} refutes the conjecture. We read
the wording in the supplied transcription:
\begin{quote}
\lq\lq Cela me fait soup\c{c}onner que la conjecture\rq\rq
\[
 X\text{ totalement }W_\infty\text{-asph\'erique}
 \Longrightarrow X\text{ filtrante}
\]
\lq\lq pourrait \^etre valable.\rq\rq
\end{quote}
We do not substitute an unproved equivalence with exactness.

For the first corrected equivalence, we checked supplied PRETOT
F.3 and TOTEX D.4 and the finite cocone criterion in
\citet[Sections 15.2(c), 15.9]{GabrielUlmer1971}.
Siftedness means
$C_{p,q}=\El_A(A(p,-)\times A(q,-))$ is nonempty and connected.
Inherited totality gives common successors in $C_{p,q}$.
For $u,v:p\to q$, a successor of
$(q;u,\id_q)$ and $(q;v,\id_q)$ gives $f,g:q\to r$
with $fu=gv$ and $f=g$. Thus $f$ post-equalizes $u,v$.
Siftedness also gives common successors in $A$, so $A$ is
filtered. The converse is the filtered-colimit criterion.

We checked the spectral-sequence steps of supplied TOTEX D.5,
D.7 and D.9; the earlier abelian one-type argument is PRETOT F.6.
For $n\geq2$, $BA$ is contractible and
$X_M=\hocolim_A K(M,n)$ is $(n-1)$-connected. Its degree-$n$
homology is $\colim M$, so the assumed truncation makes the
canonical map $X_M\to K(\colim M,n)$ a weak equivalence.
The reduced homology bar spectral sequence has first nonzero
row $M$ in degree $n$ and zero row in degree $n+1$, since
$H_{n+1}K(M,n)=0$. Hence $T_1(M)=H_{n+1}(X_M)=0$,
first proving exactness. Collapse in all degrees, applied to
$M\oplus N$, identifies the natural K\"unneth cross summand
in degree $2n$:
\[
 \colim_A(M\otimes N)\cong(\colim_A M)\otimes(\colim_A N).
\]
There is no Tor term in that degree, since two positive indices
are at least $n$. On $P_p,P_q$, this is the free abelian map
from the component set of $C_{p,q}$ to a singleton. It is an
isomorphism only if that set is a singleton. Thus $A$ is sifted
and hence filtered.

For $n=1$, common successors make the group colimit of an
abelian diagram abelian: any two representatives commute after
transport to a common object. Since $BA$ is contractible, the
fundamental-group presentation of $\hocolim_A BM$ gives
$G=\colim_A M$. The one-type assumption identifies it with
$BG$. In degree two the homology edge image is that of
$\colim_A\Lambda^2M\to\Lambda^2G$, with remaining filtration
quotient $T_1(M)$. The edge map is surjective by transporting
the two factors of an exterior product to a common successor,
so $T_1(M)=0$. After exactness is obtained, the cross summand
of $\Lambda^2(M\oplus N)$ gives the same tensor comparison
and siftedness. Conversely, filtered homotopy colimits preserve
these truncations because spheres, nullhomotopies and horns
involve finitely many simplices.

For Proposition 2.19 take $A=\Delta^{\mathrm{op}}$ and
$H=\Delta[1]/\partial\Delta[1]$. In our covariant convention,
$\El_A H$ is opposite to the usual simplex category, so its
classifying space is weakly equivalent to $|H|=S^1$.
Thus $A$ is not total. It is sifted: the cospan category in
$\Delta^{\mathrm{op}}$ on $[m],[n]$ is opposite to the simplex
category of $\Delta[m]\times\Delta[n]$, which is nonempty
and connected. This checks the variance of the counterexample.
\end{proof}

\begin{proposition}[Known, with the comparison maps audited: positive truncations]\label{D:truncation}
For every small category $A$, including empty and disconnected
categories, let $\ell:A\to\Pi A$ be groupoid completion. For
every fixed finite $n\geq1$, homotopy colimits over $A$ preserve
$n$-types if and only if
\[
 C_a=a\downarrow\ell
 =\El_A\bigl(\Hom_{\Pi A}(a,\ell(-))\bigr)
\]
is filtered for every $a$. This is supplied TOTEX D.10.
\end{proposition}
\begin{proof}[Audit of the relied-on reduction]
An object of $C_a$ is $(x,\gamma:a\to\ell(x))$; an arrow to
$(y,\delta)$ is $f:x\to y$ with $\ell(f)\gamma=\delta$.
All transition maps are bijective, and edge-path lifting identifies
$BC_a\to BA$ as the connected universal cover of its component.
For $Y:\El_A S\to\mathrm{Spaces}$, the functor
$b\mapsto\coprod_{s\in S(b)}Y(b,s)$ has the same simplicial
replacement as $\hocolim_{\El_A S}Y$. Preservation of $n$-types
therefore descends to $C_a$.

The simply connected case in TOTEX D.9 uses the reduced homology
bar complex. The constant point diagram first makes $BC_a$ an
$n$-type. For $n\geq2$, the section of
$\hocolim K(M,n)\to BC_a$ and the discreteness of the based
self-equivalence space of $K(L,n)$ trivialize its Eilenberg--MacLane
fibres over this simply connected base. Relative homology in
degree $n+1$ is zero and equals $T_1(M)$. For $n=1$, the
double-cosieve action excludes absence of common successors by
a surviving $E^2_{1,1}=\Z$ term; the abelian one-type argument
then applies. Thus $C_a$ is exact, total and, by the fixed-$n$
test, filtered. These are the specific steps checked in the
supplied proof; we do not attribute the classification to a
differently stated theorem on filtered colimits.

Conversely, filteredness makes each $BC_a$ contractible. The
pullback of $BC_a\to BA$ along $\hocolim_A X\to BA$ is
modeled by $\hocolim_{C_a}X\pi$: on the bar construction, a
lift of a chain is determined by a path class from $a$ to its
first vertex. Filtered homotopy colimits preserve $n$-types;
covering maps induce isomorphisms on homotopy groups in
degrees at least two. This proves the converse componentwise.
For empty $A$ both assertions hold vacuously.
\end{proof}

\begin{remark}[Proved: typing corrections and scope]\label{D:typing}
At II manuscript page 36, for $f:X\to Y$, finality must use
$y\backslash X$ rather than the printed $y\backslash Y$.
The latter always has initial object $\id_y$ and would make
the condition vacuous. At XII manuscript page 88,
$H'\in(U')^\wedge$, so the correctly typed isomorphism is
\[
 U'/t'\simeq U'/H',
\]
not $U'/t'\simeq U/H'$. A presheaf on $U$ is a covariant
set-valued functor on $U^{\mathrm{op}}$; its usual element
category is opposite to the corresponding covariant element
category. Opposition preserves the classifying-space homotopy type.

Definition 1 at XII manuscript pages 88--89 quantifies over
all left Kan extensions and equivalently over ordinary final
images. It is the notion used here, distinct from the
test-category diagonal condition in \emph{Pursuing Stacks}.
We inspected the supplied transcription, not an independent
copy of the authoritative edition. These are mathematical
typing corrections, not a new editorial collation.
\end{remark}

\subsection{Literature and provenance}
The access table in \texttt{TOTEX2\_report.md} records every
requested source, exact inspected theorem numbers, failed
retrievals and search terms. We obtained Isbell (1968) and
Gabriel--Ulmer (1971), which had not been opened in the
supplied record. We did not obtain the bodies of Isbell (1974),
Oberst (1968) or Ad\'amek--Reiterman (1975). No claim about
their unopened contents is used as a proof. Each bibliography
entry records body-access status.

The coefficient-ring criterion is the remark on page 296 of
\citet{IsbellMitchell1976}; the weak-terminal and bounded-hom-set
exclusions are its Corollaries 4.3--4.4, Proposition 4.5, and
Corollary 4.6. The equivalence printed in
\citet[II.3]{AdamekReiterman1976} is false.
\citet[Remark A.4]{AdamekKoubekReiterman1981} explicitly corrects
the earlier 1975 remark; their counterexample also refutes the
assertion repeated in the 1976 paper, which A.4 does not explicitly cite.
The fixed-point/exactness question in
\citet[Problem 2.5]{Husainov2002} has a weaker hypothesis
than totality and does not decide our main question.

For monoid acts, preservation of monomorphisms and of finite
limits are distinct meanings of flatness;
\citet[Definition 1.6, Proposition 1.8, Theorems 2.32 and
2.49]{HemelaerRogers2021} separates them. We also inspected
\citet[Proposition 2.4 and the paragraph after Lemma 2.6]{LaanStrongFlat}.
Neither notion may be
silently identified with flatness of the constant right $\Z A$-module.

We make no priority claim. Searches included exact colimits,
fixed-point properties, homotopy-flat actions, totally aspherical
categories, Ramsey classes, and exact titles and DOIs of the
requested sources. No opened source in these searches supplied
the missing totality-to-exactness implication or a counterexample.
This reports the search performed, not a claim that the question
is open in the mathematical literature.
